\cotefile{148}
\scanpage{44}
\mspage{1}
% archive sheet number: 43
\uline{Groupoïdes de Teichmüller.}

\textcircled{1} Cas topologique.

La variante la plus simple est
\[
\mathscr{T}_{\mathrm{top}}^{+}\quad\text{ou}\quad\mathscr{T}_{\mathrm{top}}
\]
(\uline{groupoïde de Teichmüller \uuline{orienté} topologique})

% slanted
\margin{Finalement, tant qu'on ne travaille pas au point de vue arithmétique, il semble que c'est le \unsure{gr.|gros} $\mathscr{T}^{\pm}$ qui \unsure{sera} le plus \unsure{utilisé}; \deleted{\unsure{tout}} c'est lui qui mérite donc ici la notation la plus simple $\mathscr{T}$, et le \unsure{gr.\ or.} orienté de s'appeler $\mathscr{T}^{+}$.}

\uline{Objts}: surfaces à \uline{bord} compactes orientées

\uline{flèches}: classes d'isotopie d'iso (conservant l'orientation)

Variante immédiate
\[
\mathscr{T}_{\mathrm{top}}^{\pm}
\]
(\uline{groupoïde de Teichmüller \uuline{biorienté} topologique}): mêmes objets, flèches sont les classes d'isotopie d'homéomorphismes (ne conservant pas \unsure{néc\textsuperscript{t}.|néc.} l'orientation).

Soit $[\pm1]$ le groupoïde ayant un seul objet $e$ (par exemple $e=\varnothing$) et comme groupe d'autom. $\pm1$. On a un morphisme de signature
\diagram{\begin{tikzcd}[ampersand replacement=\&]
\mathscr{T}_{\mathrm{top}}^{\pm}\arrow[r,"\chi"] \& {[\pm1]}
\end{tikzcd}}
dont la \uline{fibre} \deleted{est} (ou noyau) est $\mathscr{T}_{\mathrm{top}}^{+}=\mathscr{T}_{\mathrm{top}}$.

% slanted
\margin{\unsure{Cautèle}, \unsure{désormais borner} à des surfaces connexes !}

La variante la plus complexe est
\[
S_{*}\mathscr{T}\unsure{\mathbb{D}|D}_{\mathrm{top}}^{+}\quad\text{ou}\quad S_{*}\mathscr{T}\mathbb{D}_{\mathrm{top}},
\]
``\uline{groupoïde de Teichmüller \inserted{orienté} spécial dévissé topologique}''
